\documentstyle{article}
%\documentclass[11pt,aaspp4,flushrt]{aastex}
\def\@tightleading{0.95}
\def\baselinestretch{\@tightleading}\normalsize
\textwidth 6.4in
\textheight 8.8in
\oddsidemargin 0.0in
\evensidemargin 0.0in
\topmargin -0.3truein
\raggedbottom
\parskip 8pt
\parindent 0pt
\def\be{\begin{equation}}
\def\ee{\end{equation}}
\def\etal{{\it et al. }}
\def\kms{km~s$^{-1}$~}
\def\kmsm{km~s$^{-1}$~Mpc$^{-1}$~}
\def\Hnot{H$_\circ$}
\def\dhh{$\Delta H_\circ/H_\circ$}
\def\figwidth{1~}
\def\I{{\cal I}~}
\def\Wopt{W$_{opt}$~}
\def\Whi{W$_{21}$~}
\def\W50{W$_{50}$~}
\def\L{$\cal\char'114$}
\def\M{$\cal\char'115$}
\def\LL{{\cal\char'114}}
\def\MM{{\cal\char'115}}
\def\ts{$T_{sys}$}
\def\plotfiddle#1#2#3#4#5#6#7{{\centering \leavevmode \vbox to#2{\rule{0pt}{#2}}\special{psfile=#1 voffset=#7 hoffset=#6 vscale=#5 hscale=#4 angle=#3}}}

%\pagestyle{empty}
\begin{document}
\begin{center}
\fboxsep0.2cm
\fbox{\rule[0.0cm]{0.0cm}{0.0cm}{\parbox{16cm} 
{
\centerline{\bf ALFALFA Year 1: Midterm Report}
}}}
\end{center}

\centerline {R. Giovanelli for the ALFALFA Collaboration}
\centerline{15 July 2005}


ALFALFA, the Arecibo Legacy Fast ALFA survey of extragalactic HI, initiated
observations in February 2005. Among its scientific goals are those of delivering 
an unbiased view of the local, extragalactic Universe and providing a legacy tool 
for the astronomical community at large.

\noindent {\bf Scheduling Status.}
In the five months up to the end of June of 2005, about 400 hours of telescope
time have been scheduled, of which 367 hours were succesfully used for
observations, the rest lost to hardware and system failures or practical
''overheads'' (e.g. removal of ALFA cover). We have verified that, under
normal circumstances, the ALFALFA strategy gives a science data to observing
run duration efficiency of 97\%. Scheduling of individual observing blocks
is undertaken in a deliberate, stategic manner, reviewed almost daily based
on data quality verification. The second pass of each area of sky is 
scheduled more than 2 months after the first (so that the heliocentric
doppler shift can be exploited as a signal reality check).
For reference, see the internal A2010 observing team website:
{\it http://www.naic.edu/$\sim$a2010/galaxy\_a2010.html}.

The request submitted in October 2004
for the first survey year asked for 990 hours, or 66 observing sessions between 
07.2 and 16.7 hrs, 21.7 to 03.2 hrs, LST. This request was motivated by the
plan to complete mapping of 2 strips of ``tiles'', each 4$^\circ$ wide in 
Declination. Figure 1 shows the final status of the spring 2005 A2010 
observations. As evident in the figure, the
observations made so far cover the northern galactic cap, or the
``Spring extragalactic sky'' but with a large degree of unevenness.
We hope that time allocations towards the
end of 2005 and beginning of 2006 will allow us to complete coverage.
At the current rate of time allocation, completion of the targeted tiles
will not be possible. 

\noindent {\bf Team Constitution and Governance.}
ALFALFA is an open group, which currently includes 48 members from 11 countries
and 32 different institutions. Team expertise includes all areas deemed 
necessary for the succesful management of the survey and the achievement of
its stated science goals. 
As stated in the observing proposal, the ALFALFA team has organized its governance
by forming an Oversight Committee, which includes the P.I. (Riccardo Giovanelli),
a member from a foreign institution (Noah Brosch), a graduate student (Brian Kent), 
a member of the NAIC staff (Emmanuel Momjian) and two other US--based scientists 
(Martha Haynes and
Lyle Hoffman). A set of guidelines that will regulate the various projects'
development, within and around ALFALFA, inclduing criteria for authorship, 
has been developed and posted. Several
task groups have been formed, which contribute to the coordination of the
observations, bookkeeping, software development, data processing, follow-up
work, public access to the data, research involvement of students and public 
outreach. Particular attention to the protection of students' thesis projects
has been given in the drafting of those guidelines.

\noindent {\bf Team Activities.} 
While observing procedures have been carefully streamlined and documented ---
and observations can be and have been safely carried out remotely ---, the 
team has given strong
preference to the option of observing on site, while minimizing the use
of AO team members for that purpose. In doing so, we have appreciated NAIC's
support in waving our team's, especially students', lodging expenses. This 
practice ensures a more productive and higher
rate of utilization of the scheduled telescope time, more effective communication
with Observatory staff and shorter response times to equipment failures, more
effective training of students and faster adaptation to Observatory dynamics.
The role of ``designated observer'' has been utilized, which stimulates a
heightened sense of responsibility for the quality of the survey data and
has an invaluable impact on the training of students.

Software development is still under way, but all the fundamental pieces of
data processing software are now operational. This development has taken 
place within the IDL environment at Cornell. The software has been exported 
to and is being succesfully used by
other extragalactic HI teams, such as the NGC 2903 group. Processing of the
data to the Level I, as defined in the ALFALFA proposal (noise, bandpass 
calibrated and baselined drift scans), has been carried out for all the data
taken so far. Demonstrations of data processing at institutions other than
Cornell, such as AO, Lafayette College and Union College, have been completely
succesful, showing the easy exportability of the software.

A website is maintained, with mirror sites at Cornell 
({\it http://www.egg.astro.cornell.edu/alfalfa}) and Arecibo 
({\it http://www.naic.edu/$\sim$a2010/galaxy\_a2010.html}), where updated reports
of future scheduling and observing plans, actual observations, sky coverage,
technical memos, manuals, guidelines, science and other activities are
regularly posted. Scientific results of the survey precursor observations
are also posted: papers, data tables, catalogs and, for each source, 
cross--references to optical and other catalogs. We have made an effort to
make our results available to the community and the public in a mode
compatible with National Virtual Observatory standards. We have
obtained NSF funding to develop the necessary software and 3 of our graduate
students (Kent, Masters and Spekkens) have attended the 2004 Aspen NVO School.
Kent will return for a second year to the 2005 NVO School.
The presentation of the precursor observations data set at
{\it http://www.egg.astro.cornell.edu/precursor}, as developed by team member
B. Kent, is a preliminary demonstration of those efforts.

\noindent {\bf Team Meetings.}
In order to foster team participation, to maintain team coordination and 
proper dissemination of results and progress updates, several ALFALFA
workshops have been held, namely: 
\begin {enumerate}
\item 30 May 2005 at the Brera Observatory
downtown Milan, Italy, which saw the participation of 15 team members from
Italy, France and Spain and several other researchers interested in the 
project (taking advantage of RG's and MH's vacation in Europe); 
\item June 23-24
at Cornell in Ithaca, NY, which saw the participation of 24 researchers; 
\item 
July 6-7 at Union College in Schenectady, NY, with the participation of 
faculty from Union, Cornell,
Lafayette, Wesleyan, Colgate, U. of Puerto Rico and St. Lawrence, 2
Cornell graduate students, and 14 undergraduates from Union, Cornell,
Lafayette, Colgate and UPR. This
meeting, which focussed on research experience for undergraduates and
public outreach, was also attended by the Director of the AO Angel Ramos
Visitor Center, Dr. Jos\'e Alonso. The highlight of the meeting was
an ALFALFA observing session of 90 minutes, carried out remotely by the
students themselves and succcesfully processed and analyzed by them during
the day which followed the observations. Before the meeting, the students
had communicated by email to develop a proposal, submitted to NAIC,
for how they would use the allocated time, and remain in communication
as the data are processed and analyzed. See
{\it http://caborojo.astro.cornell.edu/alfalfalog/ualfa}. It is hoped
that the students will present a poster on their results at the Jan 2006
AAS meeting.
\end{enumerate}

\noindent{\bf Science Results.}
The observations carried out between August 2004 and January 2005, as part
of the shared--risk survey precursor effort, have been fully processed and
analyzed. Two papers have been submitted for publication to the Astronomical
Journal in May 2005, and are currently in the refereeing process. The two
papers are: 
\begin{itemize}
\item {\it The Arecibo Legacy Fast ALFA Survey: I. Science Goals, 
Survey Design and Strategy} and 
\item {\it The Arecibo Legacy Fast ALFA Survey: 
II Results of Precursor Observations}. 
\end{itemize}
Copies of the papers will be made public at the completion of the refereeing 
process, presumably by late Summer. The precursor observations were mainly
used to `shake down' the system. However, a fraction of the telescope time
yielded science grade data. Those data allowed the detection of 166 
extragalactic HI sources. Of those, 62 coincide with previously known HI sources,
while optical redshifts were available for an additional 18 galaxies; thus, 52\%
of the redshifts reported were previously unknown. Of the 166 sources, 51 are
previously unidentified objects. For comparison with the HIPASS survey,
at the sensitivity levels of that survey, {\it only about half dozen of the
166 detected precursor observation sources would have been detected by HIPASS, 
had it surveyed the same region as our precursor observations covered}.
Three of these objects had HI masses less than 10$^7$ $M_\odot$.

In the data obtained in February--March 2005, a new exciting source was found
in the general region of the Virgo cluster: an assembly of several HI clouds  
with no optical counterpart. The total HI mass is quite large, approaching
10$^9$ M$_\odot$, exceeding the masses of the neighboring optically
identified galaxies. We have confirmed the reality of the detection at
Arecibo, and through the ``rapid response'' procedure enacted at NRAO, we 
have applied for and obtained VLA C--configuration time to follow up on 
our discovery. VLA observations were made on July 11 and are being reduced.

Our team has also produced a number of technical memos, which have been
posted at the AO ALFA website, as well as in the ALFALFA website, illustrating
various aspects, concerns and results of our share of the commissioning effort. 
The Cornell Ph.D. theses of K.L. Masters (July 2005) and C.M. Springob
(Aug 2005) contain substantial work of direct relevance to ALFALFA,
including software tools for the derivation of HI mass functions, local
flow field models and survey simulations.

\noindent {\bf Utilization of Observatory Resources.}
ALFALFA has aimed to minimize the need for Observatory support. Once the 
data taking observing modes were developed and tested, no further software
development was required of Observatory staff. All the data processing
and bookkeeping software has been developed at Cornell. Four members of the
AO staff are also part of the ALFALFA team: Mikael Lerner and Phil Perillat
have provided very valuable assistance in implementing the desired features
in the observing mode, as well as in understanding the characteristics of
ALFA and as consulting sources for our data processing development. As
full members of the science team, Barbara
Catinella and Emmanuel Momjian have participated in the observations, in the
development and maintenance of web tools and as members of team working
groups but their involvement in observations has been less than that of a number
of team members. During observations, an ALFALFA ``designated
observer'' is always either at the Control Room or observing remotely. In
particular, the burden of data quality is always on an ALFALFA team member
who is also conducting the observations in real time; we have not requested 
absentee or service observing. The raw data of ALFALFA is maintained 
in storage at AO; the data rate is about 1 GB per observing hour. We have
made requests from the Computer Division for use of high end workstations  
whenever observing at AO, in order to meet the moderately high demands of
processing software. All processed data is maintained in storage at Cornell.


\noindent {\bf Recommendations for the Future.}
\begin {enumerate}
\item We appreciate the efforts of AO in allocating substantial amounts of telescope
time for this project. We wish to underscore that an effort of the magnitude
required to make ALFALFA a success requires high coordination between the 
Observartory scheduling team and the observing team. Given the observing mode
adopted, which reinforces the limitations of the telescope as a transit device,
filling in gaps in the sky coverage created by the vagaries of the start time
of each observing session, is a difficult task. That task can be helped with
``look-ahead'' scheduling, so that the observing team can plan for optimal
allocation of observing time. This is particularly important with the scheduling
of the second pass of the survey, which is ideally made at a difference of
$\sim 1/4$ or 3/4 in seasonal phase. We strongly recommend that an effort be
made to accomodate this scheduling request.

\item Given the sky coverage of the survey, separated by galactic caps, the most
intense observing seasons are January to May for the northern cap, centered at
LST of 12$^h$, and August to December for the southern cap, centered at LST of
$0^h$. A ``natural low'' in the observing thus takes place in the Summer months.
Our initial proposal was submitted on October 1st, 2004, for an observing cycle
starting in February. We would propose that the cycle be shifted by 4 months,
and that the new proposal for continuation of the survey be submitted by Feb 1, 
2006, and yearly continuation requests be submitted each year following that 
date. This would allow a seamless scheduling of the
all--important northern cap time centered at 12$^h$ LST. This proposal would
require that the currently approved proposal be extended in a ``continuing
resolution'' mode through June 1st, 2006, rather than to February 1st of that
year, and that time allocation for the February--to--June period continue to
be allocated at the same sustained rate as in 2005. Were this adjustment 
granted, the project would resume at the survey yearly cycle stated by NAIC.
\end{enumerate}

\eject
\begin{figure}[ht]
\plotfiddle{sched_05_75.eps}{2.0in}{0}{36}{36}{-360}{-50}
\plotfiddle{sched_05_60.eps}{2.0in}{0}{36}{36}{-355}{-110}
\plotfiddle{sched_05_45.eps}{2.0in}{0}{36}{36}{-360}{-182}
\vskip 2.45in
\caption[]{\small
{Actual A2010 sky coverage for Feb-July 2005. Ideally, ALFALFA
would have covered homogeneously the region extending
from $7^h30^m <$R.A.$<16^h30^m$, $+8^\circ<$ Dec.$< +16^\circ$.
By the nature of ALFALFA science, survey completion of contiguous areas
is strongly desired.}}
\end{figure}

\end{document}
